%!TEX root = ./../../thesis-trier.tex

Here, in \cref{sync_max:sec:transfer}, we transfer our notions to transformation semigroups. 
Technicalities that arise due to the possible absence of an identity transformation
in transformation semigroups are discussed in \cref{sync_max:sec:diff}
where we investigate reachability and distinguishability 
w.r.t.~to differences in semigroup and monoids.

\subsection{Transferring Notions to Transformation Semigroups}
\label{sync_max:sec:transfer} 

\index{reachable subset!in transformation semigroup}
\index{distinguishable subsets!in transformation semigroup}
\index{distinguishing!transformation}

Let $S \le_S \mathcal T_X$ be a transformation semigroup.
Similarly as for automata,  
a subset $A \subseteq X$ is said to be \emph{reachable (in $S$)} if there 
exists $t \in S$ such that $\applyfunc{X}{t} = A$.
We call two subsets $A, B \subseteq X$ \emph{distinguishable (in $S$)}
if there exists $t \in S$ (called a \emph{distinguishing transformation}) such that precisely one of the two 
subsets $\applyfunc{A}{t}$ and~$\applyfunc{B}{t}$
is a singleton subset, 
i.e.,
\[
( |\applyfunc{A}{t}| = 1 \land |\applyfunc{B}{t}| \ne 1 )
\lor 
( |\applyfunc{A}{t}| \ne 1 \land |\applyfunc{B}{t}| = 1 ).
\]

The following fact is readily seen.

\begin{proposition} \label{sync_max:prop:emptyset_dist_iff_sync}
	Every non-empty set is distinguishable from the empty set in an automaton (resp.~transformation semigroup)
	if and only if 
	the automaton is synchronizing (resp.~the transformation semigroup contains a constant map). 
\end{proposition} 

\index{completely reachable!transformation semigroup}
\index{completely distinguishable!transformation semigroup}
 
Furthermore, in analogy to automata, $S$ is called \emph{completely reachable} (resp.~\emph{completely distinguishable})
if every non-empty subset of $X$ is reachable in $S$
(resp.~if every pair of non-empty and non-singleton subsets is distinguishable).

% \todo{oben auch immer pair, und nocht two gesagt? english?}
% ~$S$ is called \emph{completely distinguishable}

In \cref{prelim:remarks_trans_semigrp_semiautomata},  
we associated automata with transformation semigroups.
By the given construction, automata-theoretical results
transfer to transformation semigroups.
For example, our results
about distinguishable pairs of subsets from \cref{sync_max:sec:distinguish} 
can be reformulated for transformation semigroups.
For reference,
we state a reformulation of \cref{sync_max:prop:dist_sync} and \cref{sync_max:thm:char_distinguishable}
for finite $X$ (where we only list, as an example,
a selection of equivalences from \cref{sync_max:thm:char_distinguishable}).

\begin{proposition}
	\label{sync_max:prop:semigrp:dist_sync}
	Let $S \le_S \mathcal T_n$
	be a completely distinguishable semigroup
	with $n \ge 3$. Then~$S$
	contains a constant map.
\end{proposition}
\begin{proof}
	Construct an automaton $\mathcal A$
	as in \cref{prelim:remarks_trans_semigrp_semiautomata}
	such that $\TS{\mathcal A} = S$.
	Then two sets are distinguishable in~$\mathcal A$
	if and only if they are distinguishable in~$S$.
	So, $\mathcal A$ is completely distinguishable and has at least three states.
	Thus, by \cref{sync_max:prop:dist_sync}, it contains a synchronizing word~$u$,
	which induces a constant map in~$\TS{\mathcal A}$.
\end{proof}

\begin{theorem}
	\label{sync_max:thm:semigrp:char_distinguishable}
	Let $S \le_S \mathcal T_n$ with $n \ge 3$.
	Then the following are equivalent:
	\begin{enumerate}	
		\item \label{dist:ts:item:all_pairs}
		$S$ is completely distinguishable.		
		
		\item \label{dist:ts:item:subsets} For every $A, B \subseteq [n]$
		with $A \subsetneq B$ and $A$ non-empty
		there exists $t \in S$ such that $1 = |\applyfunc{A}{t}| < |\applyfunc{B}{t}|$.
		
	%	\item \label{dist:ts:item:subset_single_point} For every non-empty $A \subseteq [n]$ and $p \in Q \setminus A$ there exists $t \in S$ 
	%	such that $\applyfunc{p}{t} \notin \applyfunc{A}{t}$
	%	and $|\applyfunc{A}{t}| = 1$.		
				
		\item \label{dist:ts:item:single_point} For every $p \in [n]$ there exists $t \in S$ such that
		$\applyfunc{p}{t} \notin \applyfunc{[n] \setminus \{p\}}{t}$ and
		$|\applyfunc{[n] \setminus \{p\}}{t}| = 1$.
				
		\item \label{dist:ts:item:compress} For every $A \subseteq [n]$
		with $|A| \ge 2$ and $p \in [n] \setminus A$
		there exists $t \in S$
		such that $\applyfunc{p}{t} \notin \applyfunc{A}{t}$
		and $|\applyfunc{A}{t}| < |A|$. 
		
		\item \label{dist:ts:item:dcps} For every two distinct $A, B \subseteq [n]$ 
		with $\min\{ |A|, |B|\} \ge 2$ 
		there exists $t \in S$
		such that $3 \le |\applyfunc{A}{t}| + |\applyfunc{B}{t}| < |A| + |B|$
		and $\applyfunc{A}{t} \ne \applyfunc{B}{t}$.		
	\end{enumerate}
\end{theorem}
\begin{proof}
 This follows similarly as in the proof of \cref{sync_max:prop:semigrp:dist_sync}
 by considering an automaton $\mathcal A$ with $\TS{\mathcal A} = S$.
\end{proof}


\subsection{Differences Between Monoids and Semigroups}
\label{sync_max:sec:diff} 


\paragraph{Reachability.} In a transformation semigroup $S \le_S \mathcal T_X$ every proper
subset is reachable in~$S$ if and only if it is reachable in $S \cup \{ \operatorname{id}_X \}$. So, for proper subsets, absence of an identity
element makes no difference. For finite $X$,
$X$ is reachable if and only if $S$ contains the identity transformation.

\begin{proposition}\label{sync_max:prop:reach_whole_set}
	Let $S \le_S \mathcal T_n$.
	Then $[n]$ is reachable in $S$
	if and only if $\operatorname{id}_{[n]} \in S$.
\end{proposition}
\begin{proof}
	If $\operatorname{id}_{[n]} \in S$,
	then clearly $[n]$ is reachable. 
	Conversely, suppose $\applyfunc{[n]}{t} = [n]$
	for some $t \in S$.
	Then $t$ is a permutation, and so $\operatorname{id}_{[n]} = t^i \in S$
	for some $i > 0$.
\end{proof}

As a subset is reachable in an automaton $\mathcal A$ if and only if it is reachable in~$\TM{\mathcal A}$,
we obtain the following
equivalences.

\begin{proposition}
	Let $\semiautomat{A}{Q}{\Sigma}{\delta}$.
	Then the following are equivalent.
	
	\begin{enumerate}
		\item $\mathcal A$ is completely reachable.
		
		\item $\TM{\mathcal A}$ is completely reachable.
		
		\item $\TS{\mathcal A} \cup \{\operatorname{id}_Q\}$ is completely reachable.
	\end{enumerate}	
\end{proposition}

%The identity element (resp.~the empty word) is only relevant for the reachability of the whole state set.
Most of our results about reachable subsets in automata from \cref{sync_max:sec:reachable} do not 
use the empty word. Hence they carry over to transformation
semigroups via the correspondence discussed in \cref{prelim:remarks_trans_semigrp_semiautomata}. 


%The notion of complete distinguishability neglects the empty set and pairs
%of subsets where precisely one subset is a singleton subset.
\paragraph{Distinguishability.} For automata (resp.~transformation monoids), 
the empty word (resp.~identity transformation)
distinguishes any singleton subset from any non-singleton subset.
For transformation semigroups, as the following example shows,
distinguishability of non-singleton subsets from singleton subsets is not automatically fulfilled.

\begin{example}
	\Cref{sync_max:fig:no-singleton} 
	depicts two automata where all non-empty subsets
	are indistinguishable from any singleton
	set in $\TS{\mathcal A}$.
%	Figure~\ref{prelim:fig:2-1-sets-dist} depicts an automaton~$\mathcal A$
%	where all non-empty subsets
%	are indistinguishable from any singleton
%	set in $\TS{\mathcal A}$.
%	\FIG~\ref{sync_max:fig:two-state}
%	gives a two-state automaton where no singleton subset is distinguishable
%	from the whole state set in $\TS{\mathcal A}$.
	Note that transformation semigroups over sets of size at most two
	are always completely distinguishable.

\begin{comment}
	\begin{figure}
		\centering 
		\begin{tikzpicture}
			\node[state] (1) {1};
			\node[state, right = of 1] (2) {2};
			\node[state, right = of 2] (3) {3};
			\node[state, right = of 3] (4) {4};
			\node[state, right = of 4] (5) {5};
			\node[state, right = of 5] (6) {6};
			
			%		\path[->] (1) edge [above] node {$a$} (2)
			%		(2) edge [loop above] node {$a$} (2)
			%		(3) edge [above] node {$a$} (4)
			%		(4) edge [loop above] node {$a$} (4)
			%		(5) edge [above] node {$a$} (6)
			%		(6) edge [loop above] node {$a$} (6);
			\path[->] (1) edge [above, pos=.11, out=45,in=135] node {$a$} (6)
			(2) edge [above, pos=.11, out=45,in=135] node {$a$} (6)
			(3) edge [above, pos=.11, out=45,in=135] node {$a$} (6)
			(4) edge [above, pos=.11, out=45,in=135] node {$a$} (6)
			(5) edge [above, pos=.11, out=45,in=135] node {$a$} (6)
			(6) edge [loop above] node {$a$} (6);
		\end{tikzpicture}
		\caption[Non-empty sets indistinguishable from singleton set in $\TS{\mathcal A}$]{Automaton $\mathcal A$ where every non-empty set is indistinguishable from any singleton set in $\TS{\mathcal A}$. Note that in $\TM{\mathcal A}$, every non-singleton subset is distinguishable from a singleton subset
			by the identity transformation.} 
		\label{prelim:fig:2-1-sets-dist}
	\end{figure}	
	
	% dieses beispiel doch irgendwie klar?????? tow-state immer compl dist
	% was will es zeigen? vielleicht in extending to transformation semigroups teil besser aufgehoben um nochmal zu zeigen für two-state folgt es nicht
	\begin{figure}
		\centering 
		\begin{tikzpicture}
			\node[state] (1) {$1$};
			\node[state, right = of 1] (2) {$2$};
			
			\path[->] (1) edge [bend left, above]  node {$a$} (2)
			(2) edge [loop right] node {$a$} (2)
			(2) edge [bend left, above]  node {$b$} (1)
			(1) edge [loop left] node {$b$} (1);
		\end{tikzpicture}
		\caption[Two-state automaton where whole state set indistinguishable from $1$-set.]{Automaton $\mathcal A$ 
			where$\{1,2\}$ is indistinguishable from any singleton set in $\TS{\mathcal A}$.  (but these pairs of sets are distinguishable in $\TM{\mathcal A}$).}
		\label{sync_max:fig:two-state}
	\end{figure}
\end{comment}
\begin{figure}
	\centering 
	\begin{subfigure}{0.6\textwidth}
		\centering 
		\begin{tikzpicture}[node distance = 15.5cm and 0.75cm]
			\tikzstyle{state}=[circle,draw,minimum size=24pt,inner sep=1pt,blur shadow,fill={rgb:black,0.4;white,10}]
			\node[state] (1) {1};
			\node[state, right = of 1] (2) {2};
			\node[state, right = of 2] (3) {3};
			\node[state, right = of 3] (4) {4};
			\node[state, right = of 4] (5) {5};
			\node[state, right = of 5] (6) {6};
			
			%		\path[->] (1) edge [above] node {$a$} (2)
			%		(2) edge [loop above] node {$a$} (2)
			%		(3) edge [above] node {$a$} (4)
			%		(4) edge [loop above] node {$a$} (4)
			%		(5) edge [above] node {$a$} (6)
			%		(6) edge [loop above] node {$a$} (6);
			\path[->] (1) edge [above, pos=.11, out=50,in=135] node {$a$} (6)
			(2) edge [above, pos=.11, out=50,in=135] node {$a$} (6)
			(3) edge [above, pos=.11, out=50,in=135] node {$a$} (6)
			(4) edge [above, pos=.11, out=50,in=135] node {$a$} (6)
			(5) edge [above, pos=.11, out=50,in=135] node {$a$} (6)
			(6) edge [loop above] node {$a$} (6);
		\end{tikzpicture}
		%	\caption{Automaton $\mathcal A$ where every non-empty set is indistinguishable from any singleton set in $\TS{\mathcal A}$.}
		\caption{Every non-empty set is indistinguishable from any singleton set in $\TS{\mathcal A}$.}
				\label{prelim:fig:2-1-sets-dist}
	\end{subfigure}\hfill
	\begin{subfigure}{0.36\textwidth}
		\centering
		\begin{tikzpicture}
			\tikzstyle{state}=[circle,draw,minimum size=24pt,inner sep=1pt,blur shadow,fill={rgb:black,0.4;white,10}]
			\node[state] (1) {$1$};
			\node[state, right = of 1] (2) {$2$};
			
			\path[->] (1) edge [bend left, above]  node {$a$} (2)
			(2) edge [loop right] node {$a$} (2)
			(2) edge [bend left, below]  node {$b$} (1)
			(1) edge [loop left] node {$b$} (1);
		\end{tikzpicture}
	\caption{$\{1,2\}$ is indistinguishable from any singleton set in~$\TS{\mathcal A}$.}
		%	\caption[Two-state automaton where whole state set indistinguishable from $1$-set.]{Automaton $\mathcal A$ 
			%		where$\{1,2\}$ is indistinguishable from any singleton set in $\TS{\mathcal A}$.  (but these pairs of sets are distinguishable in $\TM{\mathcal A}$).}
		\label{sync_max:fig:two-state}
	\end{subfigure}
	\caption[Non-empty sets indistinguishable from singleton set in $\TS{\mathcal A}$]{Two automata $\mathcal A$ where every non-empty set is indistinguishable from any singleton set in $\TS{\mathcal A}$. Note that in $\TM{\mathcal A}$ every non-singleton subset is distinguishable from a singleton subset by the identity transformation.} 
\label{sync_max:fig:no-singleton} 
\end{figure}
\end{example} 

Distinguishability of sets of size at least two
from any singleton subset is equivalent to the property
that for any two distinct points there exists
a transformation mapping those two points to distinct points, which is also equivalent
to distinguishability of
singletons from $2$-sets.
This is fulfilled when the transformation semigroup 
contains the identity transformation.
Next, we list other sufficient conditions
implying that for any two distinct points there exists a transformation
mapping them to distinct points.

\begin{lemma} \label{sync_max:lem:3sets-two-states-injective}
Let $S \le_S \mathcal T_n$.
Suppose for distinct $p,q \in [n]$
there exists $r \in [n]$
such that $\{p,r\}$ and $\{q,r\}$
are distinguishable.
Then for distinct $p,q \in [n]$
there exists $t \in S$ such that $\applyfunc{p}{t} \ne \applyfunc{q}{t}$.
\end{lemma} 
\begin{proof}
	Let $p,q \in [n]$.
	Then there exists $t \in S$ distinguishing $\{p, r\}$
	from $\{q,r\}$ for some $r \in Q$.
	 Without loss of generality, suppose $\applyfunc{p}{t} = \applyfunc{r}{t}$, i.e.,
	$\{p,r\}$ is mapped to a singleton set.
	Then $\applyfunc{q}{t} \ne \applyfunc{p}{t}$
	as $\{q, r\}$ is mapped to a $2$-set.		
\end{proof} 

For example, if there exists $r \in [n]$
such that, for all distinct $p,q \in [n]$,
some word $u \in \Sigma^*$
distinguishes $\{r,p\}$ from $\{r,p,q\}$,
then $u$ distinguishes $\{r,p\}$
from $\{r,q\}$. So, \cref{sync_max:lem:3sets-two-states-injective}
yields that every singleton subset is distinguishable
from every subset of size at least two.

\begin{proposition}\label{sync_max:prop:3sets-two-states-injective}
	Let $S \le_S \mathcal T_n$ and $n \ge 3$.
	Suppose pairs of distinct $2$-sets are distinguishable in~$S$.
	Then every singleton subset is distinguishable from every
	subset containing at least two elements in $S$.
\end{proposition} 
\begin{proof}
	By assumption and as $n \ge 3$,
	the condition from \cref{sync_max:lem:3sets-two-states-injective}
	is fulfilled, which gives the claim.	
\end{proof}

\begin{proposition}
	A transformation semigroup $S \le_S \mathcal T_X$
	is completely distinguishable if and only if the semigroup $S \cup \{\operatorname{id}_X\}$
	is completely distinguishable.
\end{proposition}
\begin{proof}
	Complete distinguishability does not require the distinguishability of singleton subsets from non-singleton subsets.
	Since the identity transformation distinguishes
	two subsets if and only if precisely one equals a singleton subset,
	we obtain the claim.
\end{proof}

\begin{corollary}
	Let $\semiautomat{A}{Q}{\Sigma}{\delta}$.
	Then the following are equivalent.
	
	\begin{enumerate}
		\item $\mathcal A$ is completely distinguishable.
		
		\item $\TM{\mathcal A}$ is completely distinguishable.
		
		\item $\TS{\mathcal A}$ is completely distinguishable.
	\end{enumerate} 
\end{corollary}

Besides pairs $A,B$ with $|A| = 1$ and $|B| \ge 2$,
complete distinguishability does not require distinguishability of a non-empty set from the empty set.
% nor empty set
% 
%Every non-empty set is distinguishable from the empty set if and only if 
%an automaton is synchronizing (resp.~the transformation semigroup contains a constant map).
%%The transformation semigroup $S = \{ \operatorname{id}_{[2]} \}$ is completely distinguishable
%%but $\{0,1\}$ is not distinguishable from the empty set.
%%However, for transformation semigroups on $[n]$ with $n \ne 2$ and $n \ge 1$
%%complete distinguishability also implies that the empty set is distinguishable
%%from every other set and any singleton set is distinguishable from every
%%other non-singleton subset.
Next, we show that for a transformation semigroup on $[n]$
with $n \ne 2$, complete distinguishability
implies that all pairs of sets where at least one
is not a singleton set are distinguishable.

\begin{proposition}
	A transformation semigroup $S \le_S \mathcal T_n$ with $n \ne 2$
	is completely distinguishable if and only if 
    all pairs of distinct $A, B \subseteq [n]$
	with $|A| \ne 1$ or $|B| \ne 1$
	are distinguishable.	
\end{proposition}
\begin{proof}
	For $n \le 1$ the claim is obvious, so assume $n \ge 3$ in the following.
	First, suppose $S \le_S \mathcal T_n$ is completely distinguishable.
	By \cref{sync_max:prop:3sets-two-states-injective},
	every singleton subset is distinguishable from every
	subset containing at least two points.
	By \cref{sync_max:prop:semigrp:dist_sync},
	the semigroup $S$ contains a constant map. So, by \cref{sync_max:prop:emptyset_dist_iff_sync},
	the empty set is distinguishable from every
	other set.	   
	  
    Conversely, if all pairs of distinct subsets
    where at least one subset is not a singleton set are distinguishable in $S$,
    then $S$ is clearly completely distinguishable.
\end{proof}